\section{Discussion}

The analyses point to three findings. First, work on running language models at the edge dominates the dataset. The most cited articles, the largest coupling clusters, the foundations of the co-citation network, and two of the three co-word clusters all concern efficient inference, compression, and privacy-preserving training. Second, the business-oriented stream, which covers agents, retrieval, and inventory or supply chain applications, is smaller, less connected to the rest of the field, and unstable in its exact composition. In the co-word map its cluster has the lowest centrality at the main threshold, in the coupling network only three business articles pass the citation threshold (without it they form two small groups), and in the co-citation network its few references sit inside a broad cluster of general reasoning and agent work. Third, a shift is beginning. Between 2025 and 2026, the share of articles with agent, retrieval, and small-business keywords rises, while the share with generic edge keywords falls.

The three streams share concerns such as cost, latency, privacy, and reliability, yet they build on different literatures. The inventory articles found here cite supply chain and operations management work, for example the generative AI studies of \textcite{jackson2024} and \textcite{wamba2023}, and are only weakly linked to the efficient-inference literature: in the co-citation network they fall into a broad background cluster rather than the clusters on compression and edge serving, and in the coupling network only one inventory article passes the threshold. Because the screening and the search strings favour precision (Table~\ref{tab:J}), the claim is not that bridging work does not exist, but that in this dataset it is rare.

The direction of the early shift agrees with an earlier bibliometric review of agentic AI in small, medium, and micro enterprises, which reports a sharp rise since 2022 driven by LLMs and multi-agent systems \parencite{olujimi2025}. The two studies differ in data, since this one uses OpenAlex journal articles from 2023 to 2026 and the earlier one used Scopus, Google Scholar, and grey literature from 2019 to 2025, and they also differ in scope. The agreement is therefore consistent with the earlier finding but does not confirm it. A broad bibliometric review of the LLM literature ended in early 2023 \parencite{fan2024}; the present data start in 2023 and, within the three selected streams, are dominated by edge deployment. That dominance has two possible explanations. The edge literature may simply be larger, or the terms of block B (on-device, edge device) may be more common in titles and abstracts than those of blocks A and C. The sensitivity check does not settle the question, so differences in size between the streams should not be read as a ranking of research importance. \todo{discuss why the streams diverge; this requires the author's interpretation}

\subsection{Theoretical Implications}
The maps suggest that research on LLMs for business operations has not yet drawn on the large body of work on efficient and local inference. The field is better described as three loosely connected streams than as one. Studies that join them, for example agent-based inventory support that runs on local models, would fill a gap that the data make visible. \todo{derive further claims from the labelled clusters and state no more than the data show}

\subsection{Managerial Implications}
Small businesses that consider an LLM for inventory tasks will find a small and mostly recent body of evidence in this dataset. Reported results should be treated with caution and tested on the business's own data. \todo{sharpen after reading the articles in the supply chain clusters}

\subsection{Limitations}
The study has several limitations, and each points to a direction for further work.
\begin{itemize}
  \item \textbf{One database.} The data come from OpenAlex, not Scopus or Web of Science, and metadata quality differs between the databases. For 16.8 million recent publications indexed in all three, OpenAlex has reference counts and internal coverage comparable to Web of Science and Scopus but fewer abstracts \parencite{culbert2025}, which matters for coupling and co-citation and, less, for the keyword-based maps used here. The search should be repeated in Scopus or Web of Science and the maps compared.
  \item \textbf{Incomplete reference data.} Only 282 of the 489 articles (57.7\%) have reference lists, which limits coupling and co-citation. Many of the missing lists come from journals that do not deposit references.
  \item \textbf{Partial year.} The year 2026 is incomplete, so yearly counts and trends for 2026 do not describe a full year.
  \item \textbf{Search sensitivity and breadth.} The search covers titles and abstracts only. Widening the language-model or device terms adds about five hundred records each, and these were not screened, so the baseline favours precision. One wider search of the inventory block (53 added articles, 542 in total) kept the separation of the business stream from the edge stream and gave the business stream its own co-citation and coupling clusters, but reduced the co-word map to two clusters (see the sensitivity check in the Results); widening the language-model and device terms was not tested.
  \item \textbf{Model-assisted screening.} Screening relied on language models in three passes. In a check of 20 low-confidence second-pass decisions the first author agreed with the model on 8 and rated the other 12, all excluded by the model, as relevant; a third assessment against the written criteria excluded all 12 again and added 4 of the 72 exclusions. Inclusion therefore depends on how strictly a ``main or important part'' is read, and no estimate of the overall accuracy of the screening is available.
  \item \textbf{Keywords.} Keywords come from the automatic assignment in OpenAlex, not from the authors.
  \item \textbf{Unequal blocks.} The three search blocks differ in how common their terms are, so the relative size of the streams partly reflects the search design.
  \item \textbf{Cluster stability and labels.} Co-word membership was stable but modularity was low. The co-citation partition changed with small changes in the data: four added articles merged the small supply chain group into a larger cluster, and the number of clusters ranged from three to five across runs. The coupling clusters exchanged most of their members when one article entered the map. Cluster names were chosen by the first author from model-assisted proposals; for coupling they rest on article titles only.
  \item \textbf{Metadata errors.} Three references whose titles do not fit their citing articles were removed, and other OpenAlex errors may remain. Author, institution, and country data are partly missing; country data are absent for 24\% of the articles.
\end{itemize}

\subsection{Future Research Avenues}
\todo{Make these specific enough to become study designs. Candidates: (i) a Scopus or Web of Science replication; (ii) a content analysis of the supply chain and inventory articles to test how far the results transfer to small businesses; (iii) empirical studies of agent-based LLM systems for inventory tasks that run on local or in-browser models, evaluated on real transaction data; (iv) screening of the records added by the widened search strings and a repeat of the maps.}

\section{Conclusions}
This study mapped 489 journal articles from 2023 to 2026 across three streams of LLM research: agents in business settings, models on the device or in the browser, and inventory management. Four findings emerge. Work on edge deployment dominates, since all ten most cited articles, all four coupling clusters, and the main co-cited foundations concern efficient inference, although the definition of the search blocks partly drives this result. The business-oriented stream is small and loosely connected, with the lowest centrality among the co-word clusters, only three business articles in the thresholded coupling map, and only a few co-cited references on generative AI in supply chains, which sit in a broad background cluster. An early shift is visible between 2025 and 2026 towards agents, retrieval, and small enterprises. Finally, the separation of the co-word clusters is stable but loose, while the co-citation partition is unstable and should not be interpreted in detail.

The study contributes a first joint science map of these streams, to the authors' knowledge and within the limits of searches in OpenAlex and Semantic Scholar, together with a reproducible pipeline that includes stability and sensitivity checks, and a research agenda. For practitioners, the evidence on LLMs for inventory management in small businesses is thin and recent and should be tested on their own data. The agenda has four parts: repeating the search in Scopus and Web of Science, screening the records that a wider search would add for the agent and device blocks, reading the supply chain and inventory articles in full, and building and evaluating agent-based inventory systems that run on local or in-browser models.

\section*{Acknowledgements}
The template and the method guidance follow a bibliometrics workshop given by Walton Wider at INTI International University (28 October 2026).

\textbf{Use of AI tools.} Claude \parencite{anthropic2026claude} was used through Claude Code for four tasks: screening of records on title and abstract (Claude Haiku 5.5, \texttt{claude-haiku-5-5}, 8 instances, and Claude Sonnet 5.5, \texttt{claude-sonnet-5-5}, 3 instances, both on 28 October 2026; the same Sonnet model, 2 instances, for the wider-search check, and Claude Opus 5.5, \texttt{claude-opus-5-5}, for a third assessment of 86 records, both on 29 October 2026), writing and running the data-processing scripts, proposing cluster names from the members of each cluster, and drafting and revising the text (Claude Sonnet 5.5 and Claude Opus 5.5). Table~\ref{tab:tools} gives the steps and dates. The first author wrote the research questions, search strings, screening criteria, thresholds, and decisions, rated 20 low-confidence screening decisions by hand, decided to apply the written criteria in the third assessment, chose the cluster names, and reviews and takes responsibility for the content. The model is not an author. The screening outputs, with the stated reason for each record, are kept with the data. This disclosure follows the editorial policy of the journal on the use of AI.

\section*{Author Contributions}
Aldian Yohanes: conceptualisation, methodology, data curation, formal analysis, validation, and writing (original draft), with the AI assistance stated in the Acknowledgements. Viny Christanti Mawardi: supervision and writing (review and editing). Daniel Jahja Surjawan: supervision and writing (review and editing). The second and third authors set the rule of three to five clusters per map. All authors approved the final article. \todo{the second and third authors should confirm their roles}

\section*{Conflict of Interest}
The authors declare no conflict of interest.

\section*{Data and Code Availability}
The search strings (Table~\ref{tab:search}), scripts, the screened dataset, the screening labels with reasons, and the results are held in the first author's repository (folders \texttt{analysis} and \texttt{paper}) and are available on request. OpenAlex data are open (CC0). The OpenAlex API key is not shared.

Records: OpenAlex (open index, public API, retrieved 28 October 2026; the records added for the sensitivity check, 29 October 2026). Reference lists: OpenAlex, completed with Crossref. Software: Python 3.11.7 (pandas 2.2.2, networkx 3.6.1, matplotlib 3.8.3), VOSviewer 1.6.21, XeLaTeX with biblatex 3.21 and Biber 2.21. Scopus and Web of Science were not used.
